\documentclass[autodetect-engine, ja=standard, a4paper, 10pt]{bxjsarticle}

\usepackage{graphicx}
\usepackage{amsmath, amssymb}
\usepackage{booktabs}
\usepackage{url}
\usepackage{hyperref}

\title{LuaLaTeX 論文テンプレート}
\author{著者名}
\date{\today}

\begin{document}

\maketitle

\begin{abstract}
本原稿は, LaTeX Workshop環境における LuaLaTeX 論文執筆用のテンプレートです。
日本語と英語の混在文書を美しく組み版することができます。
\end{abstract}

\section{はじめに}
LuaLaTeXは, Luaスクリプトエンジンを内蔵したTeXエンジンであり, UTF-8に標準対応しています。
従来のpLaTeXやupLaTeXと比較して, フォント指定の柔軟性や多言語対応に優れています。

\section{提案手法}
数式表現の例を示します。アインシュタインのエネルギーと質量の関係式は式(\ref{eq:einstein})の通りです。
\begin{equation}
    E = m c^2
    \label{eq:einstein}
\end{equation}

また, 参考文献の引用例を示します。LaTeXに関する文献\cite{lamport1994latex}を参照してください。

\section{実験と結果}
表\ref{tab:example}に実験条件の例を示します。

\begin{table}[htbp]
    \centering
    \caption{実験条件一覧}
    \label{tab:example}
    \begin{tabular}{ccc}
        \toprule
        項目 & パラメータ & 値 \\
        \midrule
        1 & 学習率 & 0.001 \\
        2 & バッチサイズ & 32 \\
        \bottomrule
    \end{tabular}
\end{table}

\section{おわりに}
本テンプレートを活用して論文執筆をスムーズに進めてください。

\bibliographystyle{plain}
\bibliography{ref}

\end{document}
